\section{Perfiles reproducibles y evaluación de plazos}\label{sec:implementacion-perfiles-plazos}

El catálogo de perfiles conecta la selección temporal exacta de la \cref{sec:implementacion-desencadenantes} con una configuración identificable del cálculo. Su implementación comprende valores de aplicación, revisiones persistidas en PostgreSQL y una API HTTP con consultas y comandos sobre el mismo servicio autorizado. El evaluador produce resultados en memoria; el seguimiento por expediente compone ese resultado con persistencia, responsable y atención, como se describe en la \cref{sec:implementacion-plazos-persistentes}. Los avisos requieren sus propios consumidores. Esta separación permite comprobar el cálculo y sus condiciones sin atribuir a la publicación de una configuración la validez jurídica de un vencimiento.

\subsection{Definición y ejemplos de un perfil}

Cada definición conserva título, descripción, ámbito, referencias públicas, requisito temporal, plantilla de regla, política de finalización, condiciones y ejemplos reproducibles. El ámbito puede ser global declarado o restringido a un expediente; su contenido completo queda fijo desde la primera revisión. Las referencias reutilizan los valores estructurados del catálogo de calendarios. Identifican una fuente y su localizador, pero no descargan su contenido ni certifican la interpretación que haga el despacho.

Las referencias, condiciones y ejemplos forman colecciones acotadas, con identificadores distintos y enlaces comprobados dentro de la definición. El orden declarado se conserva. Cada ejemplo contiene un tiempo con su precisión original, una cantidad ordenada opcional, un calendario completo opcional, un resultado esperado y referencias de procedencia. El constructor reproduce todos los ejemplos con el algoritmo seleccionado y exige al menos una candidata matemática. También admite ejemplos de bloqueo, útiles para expresar límites como la falta de calendario o de cantidad. Un corpus formado sólo por bloqueos se rechaza porque no demuestra que la configuración pueda calcular.

Una plantilla fija contiene la cantidad junto a la regla. Una plantilla de duración ordenada conserva la unidad y un máximo opcional, pero recibe la cantidad de la evaluación concreta. La ausencia de cantidad bloquea esa operación; el máximo no la sustituye y un exceso no se recorta. Las variantes admiten días naturales o computables, meses civiles y horas transcurridas. El desplazamiento mensual disponible busca directamente el día homólogo y bloquea su inexistencia. No utiliza una equivalencia de treinta días ni decide por sí mismo la regla aplicable a la investigación complementaria.

\subsection{Representación y revisión de la configuración}

El formato binario \texttt{DPRF1} conserva la definición completa, incluidos el orden de las colecciones, los calendarios de los ejemplos y la representación temporal original. El lector limita longitudes y conteos antes de reservar memoria, reconstruye mediante los constructores y exige que la nueva codificación coincida exactamente con la entrada. Así rechaza bytes sobrantes, etiquetas inválidas y formas que requerirían una normalización silenciosa. El instante esperado de un ejemplo conserva segundos Unix, nanosegundos y desfase original, aunque la comparación matemática de instantes pueda considerar equivalentes dos representaciones con distinto desfase.

La codificación \texttt{DINP1} conserva por separado una solicitud de insumos: expediente, selección exacta, calificación temporal, requisito, regla y referencia de calendario. Mantiene la diferencia entre ausencia y desconocimiento declarado. Su reversibilidad no implica que toda solicitud representable sea calculable: los bloqueos semánticos se resuelven después de recuperar los valores. Ninguno de estos formatos modifica los cánones de hechos, resultados de audiencia o calendarios existentes.

El catálogo asigna identidad estable, revisión positiva, estado publicado o retirado, algoritmo, huellas y autor histórico. Owner puede publicar, reemplazar y retirar; reemplazo y retiro requieren la revisión esperada y un motivo. El retiro es terminal y conserva la definición y el algoritmo anteriores. El recibo \texttt{DPTX1} vincula actor, operación, identidad, acción, revisión base, algoritmo, huella de definición y motivo. Preparar calcula una propuesta sin reservar la operación; confirmar vuelve a comprobar la base y el contenido preparado.

\subsection{Evaluación y precisión del instante final}

El evaluador recibe una definición, la selección y el material exactos, la cantidad ordenada y declaraciones explícitas de aplicabilidad. Comprueba primero la integridad de los insumos. Una condición rechazada no permite ocultar un recibo alterado, otra identidad o un calendario incoherente. Un perfil privado de otro expediente y las respuestas a condiciones duplicadas o ajenas producen errores de solicitud. Las condiciones ausentes, desconocidas o rechazadas generan bloqueos identificables, al igual que un ámbito no confirmado o una incidencia declarada pendiente.

Cuando puede extraer el tiempo e instanciar la regla, conserva el cálculo parcial y su traza incluso si otra declaración impide obtener el instante final. La estructura distingue extracción, regla opcional, aritmética opcional, bloqueos e instante final opcional. La falta de cantidad no fabrica una regla; un bloqueo de extracción no fabrica un cálculo. Las revisiones históricas y los calendarios seleccionados se mantienen separados de las cabezas observadas, sin actualización implícita.

Las horas requieren la precisión y el desfase suficientes para producir una candidata UTC. Una candidata civil puede permanecer sin hora o combinarse con un corte explícito que conserva hora hasta el segundo, desfase propio, canal, referencia y cobertura civil. El desfase del corte no se hereda del acto inicial. Si falta el corte o la candidata queda fuera de su cobertura, el evaluador conserva el bloqueo y no entrega instante final. Tampoco convierte una fecha en medianoche ni completa los segundos de una hora conocida sólo hasta el minuto. La presencia de un instante final en este resultado puro no equivale a activar una evaluación persistente ni a acreditar un efecto procesal.

\subsection{Persistencia, autorización y consulta HTTP}

Las tablas \texttt{deadline\_profiles} y \texttt{deadline\_profile\_revisions} separan raíz e historia. PostgreSQL comprueba tamaños, huellas, proyección de ámbito, recibo, secuencia y transiciones. La proyección SQL de \texttt{DPRF1} sólo recupera los campos iniciales necesarios para título y ámbito; no duplica el motor aritmético. El lector Rust y el inventario de arranque decodifican la definición completa y reproducen su corpus. Esta frontera evita presentar la aceptación del encabezado por SQL como comprobación de toda la configuración.

La confirmación usa una transacción \texttt{READ COMMITTED} y el bloqueo común de auditoría. Dentro de esa frontera revalida la identidad activa, el rol Owner, el correo capturado, la colección, la revisión esperada y, para un perfil privado, el expediente activo. La captura temporal se obtiene dentro de la transacción. Revisión, recibo, auditoría y evento de cambio se confirman juntos; un rechazo revierte sus filas. La autenticación mediante Redis y la transacción PostgreSQL siguen siendo fronteras distintas, con comprobaciones de sesión en el servicio.

La colección global excluye los perfiles privados incluso para Owner. La colección de expediente autoriza primero ese expediente y después incluye sus perfiles y los globales. Owner, Litigator y Paralegal pueden consultar; los dos últimos requieren pertenencia vigente para la colección de expediente, y Client queda denegado. El cierre conserva la consulta histórica y bloquea las mutaciones privadas. La condición global de un perfil expresa disponibilidad con ámbito declarado, no aplicabilidad universal.

La API ofrece listado ligero, detalle actual, revisión exacta, historia, preparación y confirmación de publicación, reemplazo y retiro. Las respuestas conservan colección y recibo; las consultas históricas no sustituyen la revisión solicitada. El cuerpo estructurado tiene un presupuesto propio de 16~MiB, derivado del corpus y sus calendarios, con rechazo de exceso antes de la operación. La verificación HTTP se distingue de los ensayos del adaptador y de una interfaz de usuario, todavía pendiente para estos perfiles. El contrato está en \rutarepo{docs/deadline-profiles-api.md}.

\subsection{Eventos durables y despacho de dependencias}

Las migraciones de eventos y perfiles mantienen un único registro \texttt{deadline\_source\_events}. Cada nueva revisión de resolución, notificación, resultado de audiencia, calendario o perfil emite su evento dentro de la misma transacción. La entrada conserva familia, identidad, revisión, operación y expediente cuando corresponde. Las referencias tipificadas y la comprobación de operación impiden asociarla a una revisión distinta; para un perfil global el expediente permanece ausente, sin sustituirse por un identificador artificial.

La secuencia se obtiene después del bloqueo común. Puede contener huecos por operaciones revertidas, pero no permite que una transacción confirme tardíamente un evento por detrás del cursor de otra que ya observó el bloqueo. La migración no inventa eventos anteriores a su instalación. El despachador recorre el siguiente evento realmente almacenado y expande sus dependencias en trabajos; la existencia de un evento o un trabajo no significa que su reevaluación ya haya concluido.

Las migraciones \texttt{0019} conservan trabajos y dos recorridos independientes. El recorrido de eventos pagina las cabezas actuales de plazos vinculadas a la fuente modificada; mantiene evento activo y último UUID procesado hasta completar sus candidatos. El recorrido de conciliación busca plazos V1 o con seguimiento heredado sin declarar y vuelve al inicio al terminar cada pasada. Así encuentra registros incorporados después con UUID menor, sin duplicar trabajos ni atribuirles una aceptación humana. Un evento sin candidatos también puede completarse, y un hueco de secuencia no detiene el avance.

La selección distingue familia, identidad y ámbito: una resolución puede ser fuente directa o padre de la notificación elegida; un resultado exige audiencia y expediente concordantes; un calendario mantiene ámbito global, y un perfil puede ser global o privado. Una coincidencia únicamente histórica no selecciona la cabeza actual del plazo. El cierre administrativo y la pérdida de acceso del responsable no suprimen el seguimiento técnico. Los plazos retirados quedan fuera de nuevas expansiones, mientras sus trabajos anteriores conservan su historia.

Cada página toma el bloqueo común de auditoría, verifica las capturas candidatas y confirma trabajos, cursor y auditoría en una transacción. La unicidad por evento y plazo, o por plazo y política de conciliación, conserva la asignación previa. La operación de cada trabajo queda reservada frente a envíos humanos. Una respuesta perdida permite continuar desde el cursor persistido, y un fallo de auditoría revierte toda la página. El límite de página se aplica junto con el intervalo y el filtro de trabajos ya asignados antes de materializar candidatos; acota las filas devueltas, sin prometer tiempo constante para una consulta que examine muchas raíces.

El arranque verifica tablas persistentes sin herencia ni particiones, columnas, restricciones, funciones, disparadores y permisos del rol de ejecución. El inventario recorre todas las revisiones, incluido UUID cero, y comprueba las definiciones y la existencia de los autores históricos sin exigir que sus cuentas sigan activas. El respaldo incluye tablas y secuencia; la restauración debe conservar tanto definiciones y recibos como el orden disponible para continuar emitiendo eventos.

El consumidor local autentica cada trabajo y confirma su revisión técnica, resultado sin cambios o intento fallido con la frontera descrita en la \cref{sec:implementacion-plazos-persistentes}. Despacho y consumo conservan adaptadores separados y se componen de forma secuencial en \texttt{serve}. El servicio humano y HTTP producen capturas V2 con políticas explícitas; Qadra presenta revisión, causa, autoría, observaciones y vigencia sin sustituir el cálculo histórico. Estas implementaciones tienen verificaciones focales propias, separadas de la aceptación completa contra el servidor y las bases de datos reales. La agenda conjunta y las alertas siguen fuera del alcance implementado. Se conservan cálculos en días, meses y horas, sin sustituirlos por fechas manuales ni inferir aplicabilidad jurídica de la publicación de un perfil. Las decisiones y límites se documentan en \rutarepo{docs/adr/0035-versioned-deadline-profiles-and-evaluations.md}, \rutarepo{docs/adr/0037-durable-deadline-reevaluation.md}, \rutarepo{docs/deadline-dispatch.md} y \rutarepo{docs/deadline-lifecycle.md}.
